\section{Problem}\label{sec:problem}

\subsection{Formulation}
Let $I=\{1,\dots,m\}$ be the set of candidate facilities and $J=\{1,\dots,n\}$ the set of
customers. Table~\ref{tab:notation} lists the notation. Opening facility $i$ costs $f_i\ge 0$ and
provides capacity $Q_i>0$; customer $j$ has demand $d_j>0$; serving \emph{all} of $j$'s demand
from $i$ costs $c_{ij}\ge 0$. We stress that $c_{ij}$ is a total cost per pair, not a cost per
unit: in our instances $c_{ij}=\kappa\,\mathrm{dist}_{ij}\,d_j$, and the per-unit cost
$u_{ij}=c_{ij}/d_j$ is used only to compare customers of different sizes. With $y_i=1$ if
facility $i$ opens and $x_{ij}=1$ if $i$ serves $j$, the strict single-source capacitated
facility location problem (SSCFLP) is
\begin{align}
z^\ast=\min\ & \sum_{i\in I} f_i y_i + \sum_{i\in I}\sum_{j\in J} c_{ij} x_{ij} \label{eq:obj}\\
\text{s.t.}\ & \sum_{i\in I} x_{ij} = 1 && j\in J, \label{eq:assign}\\
& \sum_{j\in J} d_j x_{ij} \le Q_i y_i && i\in I, \label{eq:cap}\\
& x_{ij} \le y_i && i\in I,\ j\in J, \label{eq:link}\\
& x_{ij}\in\{0,1\},\quad y_i \in \{0,1\} && i\in I,\ j\in J. \label{eq:int}
\end{align}
The objective \eqref{eq:obj} adds the fixed cost of the open facilities to the cost of the chosen
pairs. Constraints \eqref{eq:assign} assign each customer to exactly one facility; together with
$x_{ij}\in\{0,1\}$ this is the \emph{single-source} requirement, which forbids splitting a demand
among facilities. Constraints \eqref{eq:cap} limit the load of an open facility to its capacity
and force the load of a closed facility to zero. Constraints \eqref{eq:link} are implied by
\eqref{eq:cap} for integer solutions, but they are not implied in the linear relaxation, where
they remove fractional solutions that open a small fraction of a facility and still assign whole
customers to it. Every customer must be served: there is no penalty variable for unmet demand,
which is why we call the problem \emph{strict}.

Because of \eqref{eq:assign} and \eqref{eq:int}, a solution is fully described by an
\emph{assignment} $a:J\to I$, with $x_{ij}=1$ iff $a(j)=i$ and $y_i=1$ iff $i\in a(J)$. Writing
$\ell_i(a)=\sum_{j:\,a(j)=i} d_j$ for the load of facility $i$, the assignment is feasible iff
$\ell_i(a)\le Q_i$ for all $i$, and its cost is
\begin{equation}\label{eq:za}
z(a)=\sum_{i\in a(J)} f_i+\sum_{j\in J} c_{a(j)j}.
\end{equation}
All algorithms in this paper manipulate assignments, and every reported cost is $z(a)$ recomputed
from the original data.

\begin{table}[t]\centering\small
\caption{Notation.}\label{tab:notation}
\begin{tabular}{ll}
\toprule
Symbol & Meaning \\
\midrule
$I$, $J$; $m$, $n$ & candidate facilities and customers; their numbers \\
$f_i$, $Q_i$ & fixed cost and capacity of facility $i$ \\
$d_j$ & demand of customer $j$ \\
$c_{ij}$, $u_{ij}=c_{ij}/d_j$ & total and per-unit cost of serving $j$ from $i$ \\
$y_i$, $x_{ij}$ & opening and assignment variables \\
$a$, $\ell_i(a)$, $z(a)$ & assignment, load of $i$ under $a$, cost \eqref{eq:za} \\
$z^\ast$, $z_{\mathrm{BKS}}$ & optimal and best known cost \\
$L$ & optimal value of the strong LP relaxation \\
$v_j$, $w_i$, $s_{ij}$ & duals of \eqref{eq:assign}, \eqref{eq:cap}, \eqref{eq:link} \\
$r_i$, $\bar c_{ij}$ & reduced costs of $y_i$ and $x_{ij}$ in the strong LP \\
$U$ & cost of the incumbent (upper bound) \\
$K\subseteq I$ & set of facilities kept in a restricted problem \\
$F\subseteq J$ & set of free customers of a subproblem \\
$p_i$ & score of facility $i$ (GNN probability or LP value) \\
$T$, $\alpha$ & time budget; share of $T$ given to the expansion \\
\bottomrule
\end{tabular}
\end{table}

\subsection{The strong linear relaxation and its dual}\label{sec:lp}
Replacing \eqref{eq:int} by $0\le x_{ij}$, $0\le y_i\le 1$ gives the \emph{strong LP}, whose
optimal value we denote $L\le z^\ast$. Associate multipliers $v_j$ (free) with \eqref{eq:assign},
$w_i\ge 0$ with \eqref{eq:cap}, $s_{ij}\ge 0$ with \eqref{eq:link}, and $t_i\ge 0$ with
$y_i\le 1$. The dual is
\begin{align}
L=\max\ & \sum_{j\in J} v_j-\sum_{i\in I} t_i \label{eq:dual}\\
\text{s.t.}\ & v_j - d_j w_i - s_{ij} \le c_{ij} && i\in I,\ j\in J, \label{eq:dualx}\\
& Q_i w_i + \sum_{j\in J} s_{ij} - t_i \le f_i && i\in I, \label{eq:dualy}\\
& w_i,\ s_{ij},\ t_i \ge 0. \nonumber
\end{align}
The slacks of \eqref{eq:dualx} and \eqref{eq:dualy} at an optimal dual solution are the
\emph{reduced costs}
\begin{equation}\label{eq:rc}
\bar c_{ij} = c_{ij} - v_j + d_j w_i + s_{ij}\ \ge 0,
\qquad
r_i = f_i - Q_i w_i - \sum_{j\in J} s_{ij} + t_i\ \ge 0 .
\end{equation}
They have a direct reading. $v_j$ is the marginal cost of serving customer $j$; $w_i$ is the price
of one unit of capacity of facility $i$; $s_{ij}$ is the part of the fixed cost of $i$ that the
relaxation charges to the pair $(i,j)$. Thus $\bar c_{ij}$ is how much more expensive the pair
$(i,j)$ is than what the relaxation is willing to pay for serving $j$, once capacity and the
share of the fixed cost are priced in, and $r_i$ is the part of the fixed cost of $i$ that the
customers' willingness to pay does not cover. By complementary slackness, $\bar c_{ij}=0$ for
every pair used by the LP solution and $r_i=0$ for every facility with $0<y_i<1$.
Two uses of \eqref{eq:rc} recur in the paper: $r_i$ certifies that a pruned facility cannot
improve an incumbent (Proposition~\ref{prop:rc}), and $\bar c_{ij}$ measures how badly the
current assignment of a customer disagrees with the relaxation (the regret of
Section~\ref{sec:clns}).

A remark on a tempting shortcut: the quantity $c_{ij}-v_j$, built from the assignment duals
alone, is \emph{not} a usable regret. Since $v_j$ already embeds the capacity price and the
fixed-cost share of the facility serving $j$ in the LP, $c_{ij}-v_j=\bar c_{ij}-d_jw_i-s_{ij}$ is
non-positive on all LP-used pairs and carries no ranking information; in our first implementation
it was negative for every customer. The full reduced cost $\bar c_{ij}$ is the correct quantity.

\subsection{Why the problem is hard in practice}
The SSCFLP contains the generalized assignment problem (fix $y$) and bin packing (equal costs), so
even deciding whether a set of open facilities admits a feasible assignment is NP-complete. This
has a practical consequence for any method that reduces the set of facilities: a restricted
problem can be infeasible although total capacity exceeds total demand.

On the instance families used here (Section~\ref{sec:data}), SCIP~10 with one thread does not
prove optimality within 60~s for $m=30$ and $n=150$; the certified gap $(U-L_{\mathrm{B\&B}})/U$
after 60~s was 1.8--3.6\%, and the strong LP left $(z_{\mathrm{BKS}}-L)/z_{\mathrm{BKS}}$ of
1.6--4.4\% when fixed costs dominate. For $m=50$, $n=200$ the certified gap reached 15.6\% on one
instance.

\subsection{Capacity coupling}\label{sec:coupling}
Let $J=J_1\cup J_2$ be a partition of the customers and let $a_1$, $a_2$ be optimal assignments
of the two subproblems solved independently, each with the full capacities $Q_i$. The combined
assignment $a=a_1\cup a_2$ can fail in two ways:
\begin{align}
\text{(capacity)}\quad & \ell_i(a)=\ell_i(a_1)+\ell_i(a_2) > Q_i \quad\text{with }\ell_i(a_1),\,\ell_i(a_2)\le Q_i, \label{eq:overload}\\
\text{(fixed cost)}\quad & z(a_1)+z(a_2)-z(a)=\sum_{i\in a_1(J_1)\cap a_2(J_2)} f_i\ \ge 0.
\label{eq:double}
\end{align}
For instance, two groups that each send 60 units to a facility of capacity 100 are feasible
separately and infeasible together \eqref{eq:overload}, and the fixed cost of the shared facility
is counted twice in the sum of the subproblem objectives \eqref{eq:double}, so the subproblems
systematically overestimate the cost of sharing and under-use shared facilities. The
coordinated decomposition of Section~\ref{sec:clns} removes both effects by construction, and
Section~\ref{sec:res-decomp} measures what happens without it.
